% The Irreducible Commitment
% Part 4: Prior, Seed, and Supplied Conditioning

\documentclass[11pt,a4paper]{article}

\usepackage{amsmath,amssymb,amsthm}
\usepackage{mathtools}
\usepackage{enumitem}
\usepackage{booktabs}
\usepackage{hyperref}
\hypersetup{
    colorlinks=true,
    linkcolor=blue,
    citecolor=blue,
    urlcolor=blue,
    pdfauthor={Anonymous},
    pdftitle={The Irreducible Commitment}
}
\usepackage{cleveref}
\usepackage[margin=1in]{geometry}

\newtheorem{theorem}{Theorem}[section]
\newtheorem{lemma}[theorem]{Lemma}
\newtheorem{proposition}[theorem]{Proposition}
\newtheorem{corollary}[theorem]{Corollary}
\newtheorem{definition}[theorem]{Definition}
\newtheorem{example}[theorem]{Example}
\newtheorem{remark}[theorem]{Remark}

\newcommand{\cred}{\mathsf{cred}}
\newcommand{\Cond}{\mathsf{Cond}}
\newcommand{\Sol}{\mathsf{Sol}}
\newcommand{\negc}{\mathord{\sim}}
\newcommand{\conjc}{\otimes}
\newcommand{\disjc}{\sqcup}
\newcommand{\Real}{\mathbb{R}}
\newcommand{\Bool}{\mathbb{B}}

\title{The Irreducible Commitment\\[0.5em]
\large Prior, Seed, and Supplied Conditioning}

\author{Anonymous for blind review}
\date{June 2026}

\begin{document}

\maketitle

\begin{abstract}
Take the first measurement of an unknown quantity. Before you can update on the
reading, you must already hold a prior. Take the construction of a foundation
for mathematics. Before you can prove anything about the real line, you must
already believe a structure into which the proof runs. Neither act starts from
nothing. The opening commitment, the prior, the seed, the supplied conditioning,
cannot be derived inside the system it grounds; it can only be made explicit,
held as small as the task allows, and tested for self-consistency. This is an
old observation, the regress horn of the M\"unchhausen trilemma
\cite{albert1968} and Neurath's boat \cite{neurath1932}. Our contribution is not
a theorem about the world but a synthesis: we exhibit the same irreducible
commitment across probability, many-valued logic, and the foundations of
mathematics, and we mechanize the foundations instance in Lean so that the seed
appears as an explicit hypothesis rather than a hidden assumption.
\end{abstract}

%=============================================================================
\section{Introduction}
\label{sec:intro}
%=============================================================================

Measure an unknown length once. The instrument returns a number, and you want to
say what the length is. Bayes' rule gives the posterior from the likelihood and
the prior, but it needs the prior as input. There is no reading you can take, and
no rule you can apply to the reading, that fixes the prior for you. Pick a flat
prior and you have chosen one; pick Jeffreys' rule \cite{jeffreys1939} and you
have chosen another; refuse to choose and you cannot update at all. The first
inference forces a commitment that the data did not supply. This is not a defect
in the experiment. With more data the choice washes out under mild conditions,
which is the content of the Bernstein--von Mises theorem
\cite{vandervaart1998} and of Blackwell and Dubins on the merging of opinions
\cite{blackwell1962}. The washing out is itself conditional: Lindley's paradox
\cite{lindley1957,bartlett1957} and the inconsistency examples of Diaconis and
Freedman \cite{diaconis1986} show priors that never wash out. Either way the
prior was there first.

Now try to build mathematics from the ground. You want the real numbers, so you
want a theory in which they live. To state the axioms you already use a logic,
a notion of set or type, and a background structure rich enough to interpret
them. G\"odel's second incompleteness theorem \cite{godel1931} says that no such
theory strong enough to carry arithmetic can prove its own consistency from
inside; Gentzen's consistency proof for arithmetic \cite{gentzen1936} buys the
missing fact only by assuming transfinite induction the theory itself does not
contain. Reverse mathematics \cite{simpson2009} measures exactly how much
existence each theorem needs, and the answer is never zero. The construction
presupposes belief in a structure it cannot first justify.

These are the same situation seen twice. The measurement needs a prior; the
foundation needs a seed. In both, the opening move is a commitment that the
machinery downstream cannot generate, because the machinery runs only after the
commitment is in place. We call this the \emph{irreducible commitment}, and we
state the organizing claim plainly: no inference, justification, or construction
proceeds from nothing. Something is assumed at the start. The assumption cannot
be removed; the trilemma of Albert \cite{albert1968} closes the three escapes,
regress, circle, and arbitrary stop, and Neurath's image of rebuilding the boat
plank by plank at sea \cite{neurath1932} names the condition we are left in. The
point is older still in spirit: Quine's rejection of a foundation immune to
revision \cite{quine1951}, Wittgenstein's hinges that stand fast so that
inquiry can move \cite{wittgenstein1969}, and Carnap's distinction between
questions inside a framework and the practical choice of the framework itself
\cite{carnap1950}. We discovered none of this. We restate it.

What the commitment cannot do is hide. The honest response is not to abolish the
seed but to manage it: make it explicit, so a reader sees what was assumed;
minimize it, so the assumption is as weak as the task allows; and test it for
self-consistency, so the assumed structure does not collapse under its own
rules. Explicitness is the de Bruijn criterion in proof checking
\cite{debruijn1970}, where a small trusted kernel records every axiom a proof
leans on. Minimality is the working ideal of reverse mathematics
\cite{simpson2009} and of objective Bayesian priors \cite{berger1985,jaynes2003}.
Self-consistency is what Tarski's undefinability theorem \cite{tarski1936} and
the diagonal lemma force us to check, since a structure that can name its own
truth predicate breaks. The diagonal pattern itself is one mechanism, made
uniform by Lawvere \cite{lawvere1969} and by Yanofsky \cite{yanofsky2003}: the
liar, Russell's set, G\"odel's sentence, and Curry's paradox are instances of a
single fixed-point construction. Where a classical foundation meets such a fixed
point it must either restrict comprehension or absorb a contradiction.

The credence algebra of Parts~1--3 \cite{part1,part2} gives a setting where the
seed and its self-tests can be written down and checked. Values live in $[0,1]$
with $\negc c = 1-c$, $a \conjc b = ab$, and
$a \disjc b = a + b - ab$. Conditioning is not a connective but the external
admissibility relation $\Cond(j,e) = \{c \in [0,1] : ce = j\}$, supplied from
outside the object language. This separation is what lets a fixed point be an
equation rather than an explosion: the liar's credence $c = 1 - c$ has the unique
solution $\tfrac12$, and conditioning on impossible evidence, $e = 0$, leaves the
conditional value unconstrained instead of licensing everything. The supplied
conditioning is a seed of exactly the kind this paper studies, an input the
algebra requires and cannot manufacture. Turing's ordinal logics
\cite{turing1939} and Feferman's recursive progressions \cite{feferman1962} are
the precedent for adding such seeds in a controlled, iterated way, and
Solomonoff's universal prior \cite{solomonoff1964} is the limiting attempt to
fix the very first one.

This paper is a synthesis, not a discovery. We claim no new theorem about
priors, about provability, or about the world. We claim that one structural fact,
the irreducible opening commitment, recurs across probability, many-valued logic,
and foundations, and that it is worth seeing the three instances side by side and
in machine-checked form. The foundations instance is formalized in Lean: the seed
appears as an explicit hypothesis to a structure, the self-consistency checks
appear as theorems, and a small trusted kernel records what each derivation
assumes. Section~\ref{sec:correspondence} sets out the prior and the seed in
turn; the remaining sections develop the explicit-minimal-tested discipline and
its mechanization. The reward is modest and concrete: not a foundation without
assumptions, which the trilemma forbids, but a foundation whose assumptions are
named, bounded, and checked.

\section{The Correspondence}
\label{sec:correspondence}

The earlier sections kept three layers apart: truth values, consequence, and
conditioning.
That separation buys a dictionary.
Bayesian model selection, the trusted-base discipline of formal foundations,
and Cred conditioning all answer one question: what does a body of evidence
admit, and at what cost?
The columns name three readings of the same admissibility bookkeeping.
We collect them in Table~\ref{tab:correspondence} and explain each row.

\begin{table}[t]
\centering
\small
\begin{tabular}{p{0.27\linewidth}p{0.30\linewidth}p{0.30\linewidth}}
\toprule
Probability & Foundations & Cred conditioning \\
\midrule
prior &
foundational seed, trusted base &
supplied joint $j$ and evidence $e$ \\[4pt]
marginal likelihood, Occam factor, prior-volume cost &
trusted-base minimization, axiom strength &
admissibility relation $\Cond(j,e)$ \\[4pt]
sequential-Bayes convergence, the prior washes out &
iterated reflection, ordinal progressions &
iterated conditioning toward fixed points \\[4pt]
Cromwell's rule, a zero prior is unrecoverable &
G\"odel and Tarski limits &
conditioning on credence-zero evidence is unconstrained \\
\bottomrule
\end{tabular}
\caption{The central dictionary.  Each row reads one admissibility
question in three vocabularies.  The probability--foundations analogy is
largely folklore; the third column is this paper's explicit construction.}
\label{tab:correspondence}
\end{table}

\paragraph{Prior, seed, supplied data.}
Every column starts from something assumed before the evidence speaks.
In Bayesian inference that is the prior over hypotheses~\cite{jaynes2003}.
In formal foundations it is the trusted base: the seed axioms a proof system
takes for granted, the framework one adopts before checking
anything~\cite{carnap1950}.
In Cred it is the pair $(j,e)$ supplied to $\Cond$: a joint credence and an
evidence credence, both fixed from outside.
The point of the row is structural, not metaphorical.
Each system has a layer that is given, not derived, and the rest of the row
describes what that layer costs and how it updates.

\paragraph{Marginal likelihood, axiom cost, the admissibility relation.}
A Bayesian scores a model by its marginal likelihood, integrating the
likelihood against the prior.
Spreading prior mass over a large parameter space carries a volume penalty,
the Occam factor: a model that could have predicted anything is charged for
the prior volume it spends~\cite{jaynes2003}.
Foundations charges an analogous price.
Reverse mathematics measures a theorem by the weakest axioms that prove it,
so axiom strength becomes a currency~\cite{simpson2009}; the de Bruijn
criterion asks that a proof reduce to checking against a small trusted
kernel, minimizing what must be believed~\cite{debruijn1970}.
In Cred the corresponding object is the admissibility relation
$\Cond(j,e)=\{c\in[0,1]:ce=j\}$.
It records exactly which conditional credences the supplied data allow, and
nothing else.
The volume of $\Cond(j,e)$ is the analogue of prior cost: a singleton when
$e>0$, all of $[0,1]$ when $e=j=0$, empty when $e=0<j$.

\paragraph{Convergence, reflection, iterated conditioning.}
Run any of the three forward and ask whether it settles.
Under repeated independent observations, Bayesian posteriors concentrate and
the choice of prior washes out: the Bernstein--von Mises theorem makes this
precise~\cite{vandervaart1998}, and merging-of-opinions results show two
agents with different priors converge as evidence
accumulates~\cite{blackwell1962}.
Foundations has its own ascent.
Adding consistency or reflection statements and iterating along the ordinals
produces progressions of theories of increasing
strength~\cite{turing1939,feferman1962}.
In Cred the matching move is iterated conditioning.
Repeatedly applying $\Cond$ drives a credence toward the fixed points
of negation; the liar lands on $\tfrac12$, the
truth-teller leaves the whole interval fixed.
The first two columns are established; the third is the construction this
paper supplies.

\paragraph{The unrecoverable boundary.}
Each column has a point of no return.
Cromwell's rule warns the Bayesian never to assign prior probability zero to
a possibility, because no finite evidence can lift a zero prior back to
positive mass~\cite{jaynes2003}.
Foundations meets its own walls: G\"odel's incompleteness theorems bound what
a consistent system can prove about itself~\cite{godel1931}, and Tarski's
theorem forbids a language from defining its own truth
predicate~\cite{tarski1936}.
Cred's boundary is the zero-evidence case.
When $e=0$ the chain-rule equation $ce=j$ imposes no constraint on $c$, so
$\Cond(0,0)=[0,1]$ and conditioning on impossible evidence is unconstrained.
The difference is the design choice of the whole program.
A zero prior is a trap because probability conditions by division; Cred
conditions by the admissibility relation instead, so credence-zero evidence
yields no value rather than a contradiction.
There is no ex falso here, by construction.

\paragraph{What is folklore and what is new.}
The probability--foundations pairing in the first two columns is not our
discovery.
The link between Occam factors and description length, between Bayesian
convergence and the washing out of priors, and between inductive inference
and formal provability has been drawn before, from algorithmic information
theory to the philosophy of induction.
Our contribution is the third column and the claim that it lines up row by
row with the other two.
Cred conditioning is the external relation $\Cond$, not an object-language
connective; that externality is what lets the zero-evidence row read as
\emph{unconstrained} rather than \emph{explosive}.
We state the correspondence as an organizing analogy, not a theorem: no
functor or bi-interpretation is claimed between the columns.
What is proved is the Cred column, in Lean, and its internal coherence with
the fixed-point and no-explosion results of the preceding sections.

\section{A Minimal Model}
\label{sec:toymodel}

A single iterated map on $[0,1]$ exhibits all four features of the
framework before any logic enters.
Fix a target credence $t \in [0,1]$ and define one refinement step:
\[
  f_t(x) = \frac{x + t}{2}.
\]
Starting from a seed $x_0$, the orbit is $x_n = f_t^n(x_0)$.

\begin{example}[Toy model]
\label{ex:toymodel}
Take $t = \tfrac{1}{2}$ (the liar fixed point of $\negc$).
Seed $x_0 = 0$ gives $x_1 = \tfrac{1}{4}$.
Seed $x_0 = 1$ gives $x_1 = \tfrac{3}{4}$.
The two orbits differ from step one.
\end{example}

Four properties hold; each is a machine-checked lemma in
\texttt{Cred/ToyModel.lean}.

\begin{theorem}[P1: Seed matters locally]
\label{thm:p1}
Distinct seeds produce distinct finite orbits.
Formally, there exist seeds $x_0, x_0' \in [0,1]$ and a step $n \in \mathbb{N}$
such that $f_t^n(x_0) \neq f_t^n(x_0')$.

\emph{Lean:} \texttt{toyA\_seed\_matters}.
For $t = \tfrac{1}{2}$, seeds $0$ and $1$ already differ at $n = 1$.
\end{theorem}

The seed is unavoidable: no finite prefix of the orbit erases the starting
commitment.
This is the irreducibility claim at finite time.

\begin{theorem}[P2: Seed washes out]
\label{thm:p2}
The signed distance to the target satisfies
\[
  f_t^n(x_0) - t = \Bigl(\tfrac{1}{2}\Bigr)^n (x_0 - t),
\]
so $|f_t^n(x_0) - t| = (\tfrac{1}{2})^n |x_0 - t|$,
and $f_t^n(x_0) \to t$ as $n \to \infty$ from any seed $x_0$.

\emph{Lean:} \texttt{toyA\_contraction\_bound} (the geometric bound)
and \texttt{toyA\_contraction\_tendsto} (the \texttt{Tendsto} statement).
\end{theorem}

The geometric decay is the local analogue of merging of opinions:
asymptotically the prior is forgotten and only the target survives
\cite{blackwell1962,vandervaart1998}.
The bound is explicit and the limit is a genuine \texttt{Tendsto} statement
in Mathlib, not just an $\epsilon$-$\delta$ argument.

\begin{theorem}[P3: Zero seed is absorbing]
\label{thm:p3}
Under the product map $g_e(x) = x \conjc e$, the zero credence is
absorbing: $(g_e)^n(0) = 0$ for all $n$.
A seed of $0$ never recovers to a positive credence.

\emph{Lean:} \texttt{toyA\_absorbing\_zero}.
The companion lemma \texttt{toyA\_absorbing\_zero\_no\_ex\_falso}
witnesses each iterate as a \texttt{Conditioning 0 0} value,
connecting absorption directly to \texttt{conditioning\_zero\_any}.
\end{theorem}

This is the iterated face of Cromwell's rule \cite{jaynes2003}: impossible
evidence (credence $0$) provides no constraint on belief and stays impossible
under repeated conjunction.
The model enforces no ex falso: the zero orbit never yields a
non-zero credence, no matter what $e$ is applied.

\begin{theorem}[P4: Limit satisfies the fixed-point equation]
\label{thm:p4}
The target $t$ is a genuine fixed point of $f_t$: $f_t(t) = t$.
The orbit from any seed converges to $t$ (P2), and $t$ satisfies
the self-consistency equation the base level converged toward.

\emph{Lean:} \texttt{toyA\_target\_fixed} (the fixed-point equation)
and \texttt{toyA\_limit\_fixed} (the conjunction of P2 and the fixed-point
equation as a single statement).
\end{theorem}

The fixed-point check is one level up: the base-level iteration
converges to $t$, and $t$ satisfies $f_t(t) = t$, so the meta level
certifies the value the base level assumed.
The model is self-consistent without any external grounding.

\begin{remark}
The four properties are independent by construction.
P1 and P2 together say that the seed matters at every finite stage but
not in the limit: local irreducibility coexists with global convergence.
P3 isolates the boundary behavior that would otherwise cause explosion in a
conditional logic.
P4 closes the loop between levels.
All lemmas are proved in Lean~4 with Mathlib and carry zero \texttt{sorry}.
\end{remark}

\section{Iteration: Sequential Evidence and Iterated Reflection}
\label{sec:iteration}

Two procedures repeat a single step and approach a limit. Sequential
conditioning repeats Bayesian update. Iterated reflection repeats adding the
soundness of a theory to that theory. The procedures are not the same. They
share a shape worth stating, and they break at different boundaries.

\subsection{Updating forgets the prior}

Start with one parameter and a stream of observations. The prior is a
credence over the parameter. Each new datum multiplies it through the chain
rule and renormalizes. Run this on independent, identically distributed data
and the posterior concentrates on the truth.

Two theorems make the convergence precise. Doob's consistency result,
sharpened by the Bernstein-von Mises theorem, says the posterior is
asymptotically normal around the truth and shrinks at rate $1/\sqrt{n}$
\cite{vandervaart1998}. The proper prior survives in the
$1/\sqrt{n}$ correction and then washes out. A second result speaks to
disagreement between observers. If two agents assign positive credence to the
same events, their posteriors merge as evidence accumulates
\cite{blackwell1962}. Different starting credences converge to the same limit.

The slogan is that data forget the prior. The prior is not erased at any
finite stage; its weight decays. In the limit only the likelihood speaks.

\subsection{Reflection adds soundness}

Now replace data by proofs about proofs. A consistent theory $T$ strong enough
for arithmetic cannot prove its own consistency \cite{godel1931}. It can,
however, be extended. Add to $T$ the statement that $T$ is sound, or the
reflection principle ``if $T$ proves $\varphi$ then $\varphi$.'' Call the
result $T_1$. The new theory proves the consistency $T$ could not.

Iterate. Each $T_{n+1}$ adds the soundness of $T_n$ on top of $T_n$.
Continue past the finite stages into the transfinite, indexed by ordinals,
taking unions at limit ordinals. Turing studied the first ordinals of such a
progression \cite{turing1939}; Feferman developed the transfinite recursive
progression in full \cite{feferman1962}. Each step is sound if the previous
theory was. The sequence climbs toward a limit theory that settles more
arithmetic than any finite extension.

\subsection{Where the analogy holds}

Both procedures iterate one operator and read off a limit. In updating the
operator is conditioning on the next datum; in reflection it is adjoining the
soundness of the current level. Both limits are stronger than every finite
stage and licensed by the stages below. Neither limit is reached at a finite
step. Both formalize the same intuition: standing where you are, you can see
that your present position is sound, and saying so moves you forward.

\subsection{Where it breaks}

The convergence is a theorem with hypotheses, and the hypotheses fail. Doob's
theorem secures consistency only outside a prior-null set; the null set can be
large. In infinite-dimensional models a prior can place that null set on the
truth, and the posterior then concentrates on the wrong value no matter how
much data arrive \cite{diaconis1986}. Merging needs both agents to grant
positive credence to the truth; a dogmatic prior that excludes it never
merges. Convergence is contingent on the prior, not automatic.

Reflection breaks against a different wall. Each added soundness statement
is itself an arithmetic sentence, so the new theory is again incomplete by the
same theorem that forced the first extension \cite{godel1931}. The
progression never closes into a complete theory; the incompleteness boundary
moves with it. Worse, the route through the ordinals depends on how each
ordinal is named, and the naming smuggles in strength. The limit is not a
single canonical theory but a family indexed by notations
\cite{feferman1962}.

So the analogy holds in shape and fails in kind. One iteration forgets a
finite-weight prior under a measure-theoretic condition that can fail in large
models. The other adds genuine deductive strength under a syntactic condition
that can never terminate. The first has a limit theorem with a sharp
counterexample; the second has a limit construction with no completion. Reading
them together sets a frame for the conditioning boundary of
Section~\ref{sec:boundaries}: convergence and consistency are earned
under stated hypotheses, never by fiat, and both run out exactly where their
hypotheses do.

\section{Two Boundaries}
\label{sec:boundaries}

The credence algebra meets probability and statistics at two boundaries where
a multiplicative update degenerates. One sits at evidence value zero. The
other sits where a prior loses width. Both are facts about the limit of a
product, not separate pathologies. This section states each as such and marks
which ties are this paper's sharpenings rather than prior results.

\subsection{The Unrecoverable Zero}
\label{sec:zero}

Start with a prior credence of zero on a hypothesis $A$. No finite evidence
moves it. Cromwell's rule names this: a probability set to zero or one is
never revised by conditioning~\cite{jaynes2003}. The mechanism is
multiplication. Update multiplies the prior by a likelihood ratio, and zero
times anything is zero.

The same arithmetic governs the evidence side. The chain rule is
\[
  \cred(A\mid B)\conjc\cred(B)=\cred(A\wedge B).
\]
When $\cred(B)=0$ the equation reads $\cred(A\mid B)\conjc 0=0$. Every value
in $[0,1]$ satisfies it. Conditioning on impossible evidence is unconstrained:
\[
  \Cond(0,0)=\{c\in[0,1]:c\conjc 0=0\}=[0,1].
\]
Zero is absorbing under $\conjc$, so the boundary works in both directions. A
zero prior cannot be moved. Zero evidence cannot move anything. One absorbing
element produces both readings.

\begin{remark}
The two faces are dual. A zero in the multiplicand fixes the product at zero
regardless of the multiplier; a zero in the position that update would divide
by leaves the quotient free. Division would force a choice on the null set, by
convention or by a limit. The chain rule, taken as primitive, declines to
choose. The admissible set is the whole interval.
\end{remark}

This is why impossible evidence neither forces a conclusion nor can be forced
to one. No ex falso enters. From $A$ and $\negc A$ the calculus derives no
unrelated positive conclusion; Lean records this as
\texttt{labelled\_no\_ex\_falso}. The zero-conditioning fact is
\texttt{cond\_zero\_zero\_univ}, with the underlying value-level statement
\texttt{conditioning\_zero\_any}. The contrast with a residuated implication is
sharp: an implication connective would let the impossible antecedent discharge
any consequent. The external conditioning judgment $c\in\Cond(j,e)$ keeps the
antecedent outside the formula language, so the discharge never happens.

The result is conservative. We did not discover that zero is absorbing or that
Cromwell's rule holds. The sharpening is to read both as one boundary of the
multiplicative update and to certify the no-explosion consequence in Lean.

\subsection{The Measure-Zero Verdict}
\label{sec:occam}

The second boundary is where a prior loses width. Compare two models by their
marginal likelihoods. For a model with parameter $\theta$, prior $\pi$, and
data $x$,
\[
  m(x)=\int p(x\mid\theta)\,\pi(\theta)\,d\theta
  \;\approx\; p(x\mid\hat\theta)\cdot\Omega,
\]
where $\hat\theta$ maximizes the likelihood and $\Omega$ is the Occam factor:
how much of its prior width the fitted parameter actually uses. The Bayes
factor multiplies the maximized likelihood ratio by the ratio of Occam
factors.

Width controls $\Omega$. Widen the prior on a parameter that the data pin down
and $\Omega$ shrinks, penalizing the larger model. Send the prior width to
infinity, the improper limit, and $\Omega\to 0$ at a rate the prior does not
fix. The Bayes factor then has no determinate value. Bartlett stated this
against a flat prior: the answer depends on a width that the improper prior
refuses to name~\cite{bartlett1957}. Jeffreys had already built the point into
his test design, placing a proper prior of fixed scale on the alternative so
that $\Omega$ stays defined~\cite{jeffreys1939}.

Lindley's paradox is the same factor read along a trajectory~\cite{lindley1957}.
Fix a test statistic at a borderline value, say $Z$ two standard errors from
the null, and let the sample size $n$ grow. The frequentist verdict is
constant: reject at the chosen level. The Bayesian verdict reverses, because
$\Omega$ falls like $1/\sqrt{n}$ as the posterior concentrates while the
likelihood ratio stays put. The two verdicts agree only on a fine-tuned path
that holds $Z$ fixed as $n$ grows. That path is a set of measure zero in the
data. Off it, with $Z$ free to track the data, the divergence does not arise.

The frequentist test is a Bayes factor with the Occam factor set to one. A $Z$
or $\chi^2$ test rejects on the maximized likelihood ratio,
\[
  \exp\!\big(Z^2/2\big),
\]
which is exactly the Bayes factor's first term with $\Omega=1$. The test keeps
the model's best fit and drops the width penalty. Lindley's reversal is the
gap between $\Omega=1$ and the $\Omega$ that a proper prior supplies. Naming
that implicit $\Omega=1$ is this paper's sharpening; the likelihood-ratio form
of the test statistic is standard.

Under uncertainty about the prior, the honest object is not a number. It is an
interval. Range the prior over a class and the Bayes factor ranges over the
induced set of values. Imprecise probability takes this set as primitive
rather than selecting one prior~\cite{walley1991}; robust Bayesian analysis
reads the same set as sensitivity to the prior
class~\cite{berger1985}. The interval of Bayes factors is the statistical
image of set-valued admissible conditioning. Where evidence underdetermines
the update, $\Cond$ returns a set; where the prior class underdetermines the
comparison, the Bayes factor returns a set. The two boundaries share the form:
a multiplicative update degenerates, and the admissible answer is an interval,
not a point.

\begin{remark}
What we claim is narrow. The Occam factor, Bartlett's improper-prior
indeterminacy, Lindley's paradox, and the imprecise-probability interval are
prior results, cited as such. The sharpenings are three ties: the implicit
$\Omega=1$ inside the frequentist test, the measure-zero character of the
agreement trajectory, and the match between the Bayes-factor interval and the
set-valued $\Cond$. None of these requires a new statistical mechanism. Each
reads an existing one as a boundary of the same product.
\end{remark}

\section{Relation to Existing Syntheses}
\label{sec:relation}

Two unifications already cover most of the ground this paper stands on.
Each takes a face of the problem and shows it to be one thing.
Neither covers the join we make here, and we claim only the join.

\subsection{Self-reference as one diagonal fact}

Lawvere proved that Cantor's theorem, Russell's paradox, the liar, G\"odel's
first incompleteness theorem, Tarski's undefinability of truth, and Curry's
paradox are instances of a single fixed-point lemma in a cartesian closed
category~\cite{lawvere1969,godel1931,tarski1936}.
A point-surjective map $A\to B^A$ forces every endomap of $B$ to have a fixed
point; when $B$ carries a fixed-point-free map, no such surjection exists.
Yanofsky restated this without category theory, as one diagonal argument that
generates each paradox by choosing the objects and the map~\cite{yanofsky2003}.
The synthesis is structural.
It says these phenomena share a shape.
It does not assign the liar a value or say what a theory should do once the
shape appears.

Our solution-set reading sits downstream of that lemma.
The negation map on $[0,1]$ has the unique fixed point $\tfrac12$, so the
liar and Russell equations have solution $\{\tfrac12\}$; the identity map has
every point fixed, so the truth-teller has solution $[0,1]$.
Lawvere guarantees a fixed point exists.
The credence algebra says which one, and how many, on a concrete carrier.

\subsection{Probability as the logic of plausible inference}

Cox derived the sum and product rules from consistency requirements on a
real-valued degree of belief~\cite{cox1946}, and Jaynes built probability
theory as the unique consistent extension of Boolean logic to graded
plausibility~\cite{jaynes2003}.
This unifies the other face: deductive logic is the certainty limit of one
calculus of belief, not a separate subject.
The product rule $\cred(A\wedge B)=\cred(A\mid B)\,\cred(B)$ is the chain rule
we take as primitive.

Cox and Jaynes fix the conditional by division, $\cred(A\mid B)=\cred(A\wedge
B)/\cred(B)$, and leave the case $\cred(B)=0$ undefined or excluded by a
regularity convention.
That convention is exactly where explosion re-enters: a vacuously true
conditional on impossible evidence licenses any conclusion.

\subsection{What remains open, and what we add}

Neither synthesis settles the boundary.
Lawvere classifies the paradox; Cox and Jaynes classify the inference; the
seam between them, conditioning on evidence that has been ruled out, is left
to convention on each side.
This paper works that seam.

The contribution is a three-way correspondence and its mechanization.
The same admissible set
\[
  \Cond(j,e)=\{c\in[0,1]:c\conjc e=j\}
\]
reads three ways.
As a prior, $c$ is the credence updated by positive evidence $e$ toward joint
$j$.
As a seed, it is the starting value a self-referential equation must satisfy.
As a supplied joint, it is the dependence a valuation must fix before a
conjunction has a value.
For $e>0$ the set is the singleton $\{j/e\}$; for $e=j=0$ it is all of
$[0,1]$; for $e=0<j$ it is empty.

Two classical results sharpen the boundary.
The Cromwell rule, that a prior of $0$ or $1$ can never move, is here the
no-ex-falso theorem: impossible evidence imposes no constraint, so it cannot
force a conclusion.
Lindley's paradox, where a sharp null and a diffuse alternative pull a verdict
apart~\cite{lindley1957,bartlett1957,jeffreys1939}, is here the move from a
single posterior to the admissible set $\Cond(j,e)$: when the chain rule does
not pin a value, the honest answer is the whole interval, not an arbitrary
point. This matches the imprecise-probability stance on null
events~\cite{walley1991}.

The components are not ours.
The fixed-point lemma is Lawvere's and Yanofsky's; the inference calculus is
Cox's and Jaynes's; the admissible-set boundary draws on primitive
conditional probability and coherence.
What is new is the join: one carrier on which the diagonal lemma and the
chain rule meet, with the zero-evidence boundary made explicit rather than
conventional, and the whole separation checked in Lean.
Every theorem named in this paper type-checks against the Lean kernel with no
\texttt{sorry}.
The kernel plays de Bruijn's role: a small trusted core re-checks each
proof, so the reader trusts the checker, not the author~\cite{debruijn1970}.

\subsection{Two further relocations}

Solomonoff's universal prior makes Occam's razor a prior: shorter programs get
more weight, so simpler hypotheses start ahead~\cite{solomonoff1964}.
The universality has a price.
The prior is universal only up to a multiplicative constant that depends on
the chosen reference machine, and that constant is not bounded across machines.
This is the same pattern we keep flagging.
A modelling choice that looks canonical (a universal prior, a regular
conditional on a null set, a value for an independent statement) hides a free
parameter.
The credence algebra forces the parameter into the open, as the supplied
joint or as the admissible interval.

Gentzen proved the consistency of arithmetic by transfinite induction up to
$\varepsilon_0$, trading one unprovable statement for a stronger principle the
original theory could not reach~\cite{gentzen1936}.
Trust is relocated, not eliminated; this is the regress that the
M\"unchhausen trilemma names in general~\cite{albert1968}.
Our mechanization relocates trust the same way.
It does not certify the credence algebra from nothing.
It moves trust from informal argument to the Lean kernel and the mathlib
axioms it rests on.
We claim the relocation is to a smaller, public, re-runnable core, not that it
is to no core at all.

\section{Conclusion}
\label{sec:conclusion}

The product De Morgan triplet on $[0,1]$ (product conjunction,
probabilistic-sum disjunction, standard complement) sits at an
intersection that is easy to miss: it satisfies every axiom shared by
Bayesian probability~\cite{cox1961,jaynes2003} and product
logic~\cite{hajek1998}, yet commits to neither convention at the one
point where they diverge.
At positive evidence, chain-rule conditioning agrees with Bayesian
division and with product residuation; at zero evidence, it imposes no
value.
The annihilator identity $c \cdot 0 = 0$ is the whole story: it forces
no posterior when evidence is impossible, provides the paraconsistency
witness ($v(A) = \tfrac{1}{2}$ satisfies both $A$ and $\negc A$ as
positive without constraining $B$), and explains why the collapse
homomorphism cannot extend to conditioning, since collapsed values discard
the ratio information that the conditional requires.

What this paper contributed, beyond collecting known components, is
fourfold.
First, a dictionary: LP consequence on the Kleene lattice is
\emph{exactly} positivity consequence on $[0,1]$, and K$_3$ is exactly
certainty consequence (the Lean bridge theorems \texttt{lp\_formula\_bridge}
and \texttt{k3\_formula\_bridge}).
An LP logician is implicitly doing positivity reasoning over $[0,1]$;
a K$_3$ logician is implicitly doing certainty reasoning.
The collapse is the Rosetta stone, and it works in both directions.
Second, a machine-checked toy model: every result is verified in
Lean~4 with Mathlib and zero \texttt{sorry}, meeting the de Bruijn
criterion~\cite{debruijn1970} that a claimed proof can be checked by
an independent verifier.
Third, two boundary sharpenings of the conditional impossibility
(the Lean results \texttt{cond\_bridge\_boundary} and \texttt{cond\_bridge\_fails\_interior}): the bridge
fails precisely when both conclusion and evidence are interior and the
conclusion is less credible than the evidence; it holds whenever at
least one is at a boundary.
Fourth, an honest positioning: the paper does not claim a
reconciliation of probability and many-valued logic, only a shared
algebra and an explicit account of where they part company.

The practical reading is this.
Any framework that handles conditioning must commit to something at
evidence zero: Bayesian probability leaves the conditional undefined,
product residuation fills it with $1$ (ex falso), and chain-rule
conditioning leaves it unconstrained.
None of these is the neutral option.
A system's worth is not that it has no commitment at this point, but
that it makes its one unavoidable commitment explicit, minimal, and
checkable.
The product De Morgan triplet with chain-rule conditioning makes that
commitment as small as arithmetic allows: $c \cdot 0 = 0$, nothing
more.

\appendix

\section{Priors as seeds: informative, reference, and model volume}
\label{sec:priorseed}

The prior is the Bayesian name for the seed. Before the first datum, an inference
already holds a credence over hypotheses, and no reading fixes it from inside.
Part~4 argued this in the abstract; here we read the prior as one concrete
instance of the irreducible commitment, and we follow it into model comparison,
where the seed acquires a price.

\paragraph{No prior is neutral.}
An informative prior commits a lot. It names where the parameter is likely to
sit and how tightly, and it pays for a wrong guess with slow convergence. A
reference prior commits less: Jeffreys' rule, the flat prior, the maximum-entropy
prior under stated constraints all aim to let the data speak with least
interference~\cite{jaynes2003}. Less is not none. A flat prior on a parameter is
not flat after reparametrization; maximum entropy fixes a measure and a
constraint set; Jeffreys' rule fixes an invariance. Each choice is a commitment
dressed as the absence of one. The reference prior is the seed held as small as
the task allows, never the seed removed.

\paragraph{Prior volume is the Occam factor.}
The commitment becomes a number when two models are compared. Section~\ref{sec:occam}
wrote the marginal likelihood as $m(x)\approx p(x\mid\hat\theta)\cdot\Omega$, with
$\Omega$ the Occam factor: the fraction of its prior width the fitted parameter
uses. A broad prior spreads credence over a large parameter space, so $\Omega$ is
small, and the broad model is charged for the volume it spends~\cite{lindley1957}.
A model with a diffuse prior is therefore a \emph{larger} model, not a neutral
one. Widen the prior toward the improper limit and $\Omega\to0$ at a rate the
prior never names, so the Bayes factor loses a determinate value; Bartlett read
this off the flat prior directly~\cite{bartlett1957}. The seed that looked least
committal turns out to commit the most, by claiming the most volume.

This is the same accounting as $\Cond$. The admissible set
$\Cond(j,e)=\{c\in[0,1]:c\conjc e=j\}$ records what the supplied data leave open;
its width is the analogue of prior volume. A broad prior buys a large $\Omega$
penalty exactly as a degenerate $e$ buys a wide $\Cond$. Width is the currency in
both columns, and width is never free.

\paragraph{The reference-prior dream and its machine.}
Push the reference prior to its limit and you reach Solomonoff's universal prior:
weight each hypothesis by the length of the shortest program that generates it, so
simpler hypotheses start ahead and Occam's razor becomes the seed
itself~\cite{solomonoff1964}. This is the cleanest attempt to fix the very first
commitment once and for all. The price is the same one the flat prior pays in a
cruder form. Universality holds only up to a multiplicative constant set by the
chosen reference machine, and that constant is unbounded across machines. The
seed is minimal in description length and still machine-dependent. Even the prior
built to be canonical carries a free parameter, relocated from the shape of a
distribution to the choice of a universal machine.

The lesson tracks the Part~4 thesis without adding to it. Every prior is a seed
the inference cannot derive from its data. Making it informative commits more;
making it diffuse commits more volume, not less; making it universal commits to a
machine. The honest move is not to find the neutral prior, which does not exist,
but to name the seed, hold it small, and let the comparison expose what its width
costs.

\section{A worked example: does gravity exist?}
\label{sec:gravity}

Everything above runs on one experiment. We measure a downward acceleration and
ask whether a gravitational effect exists at all. The two boundaries of
Section~\ref{sec:boundaries} (the Occam factor and the washing out of a proper
prior) and the coherence theorem of Section~\ref{sec:iteration} all show up in
explicit numbers.

Set up two models for a single reading $x$ of the acceleration, in units where
the noise has standard deviation $\sigma = 1$.
\begin{itemize}
\item $M_0$: no effect. The acceleration is exactly $g = 0$, so
$x \sim \mathcal{N}(0,1)$.
\item $M_1$: an effect of unknown size. The acceleration is some
$g \in [0, W]$, with a flat prior of width $W$, and $x \sim \mathcal{N}(g,1)$.
\end{itemize}
The width $W$ is the prior on $M_1$. It is the seed of Section~\ref{sec:correspondence}:
chosen before the reading, not read off it.

The marginal likelihood of $M_1$ factors the way Section~\ref{sec:occam}
describes,
\[
  m_1(x) = \int_0^W \frac{1}{W}\, \frac{1}{\sqrt{2\pi}}\,
  e^{-(x-g)^2/2}\, dg
  \;\approx\; \underbrace{\frac{1}{\sqrt{2\pi}}\, e^{-(x-\hat g)^2/2}}_{\text{best fit}}
  \cdot \underbrace{\frac{\sqrt{2\pi}}{W}}_{\Omega},
\]
with $\hat g = x$ when $x$ lands inside $[0,W]$. The Occam factor
$\Omega = \sqrt{2\pi}/W$ is the fraction of its prior width that the fitted
parameter uses. The Bayes factor is the maximized likelihood ratio times that
factor:
\[
  B_{10} = \frac{m_1(x)}{m_0(x)}
  \;\approx\; \underbrace{e^{x^2/2}}_{\text{likelihood ratio}}
  \cdot \underbrace{\frac{\sqrt{2\pi}}{W}}_{\Omega}.
\]

\begin{example}[A broad prior loses on a datum that favors gravity]
\label{ex:gravity-lindley}
Read $x = 2$: two standard errors from zero, so the datum favors $g \neq 0$.
The likelihood ratio is $e^{2} \approx 7.39$, decisive against $M_0$ on its
own. Now feel the width.
\begin{itemize}
\item Narrow prior, $W = 5$: $\Omega = \sqrt{2\pi}/5 \approx 0.50$, so
$B_{10} \approx 7.39 \cdot 0.50 \approx 3.7$. Evidence for gravity.
\item Broad prior, $W = 100$: $\Omega \approx 0.025$, so
$B_{10} \approx 7.39 \cdot 0.025 \approx 0.18$. Evidence \emph{against}
gravity, with the same datum.
\end{itemize}
The reading did not change. The seed did. A prior too broad on $g$ charges
$M_1$ for parameter space it never used, and $M_1$ loses even though the datum
favors $g \neq 0$. This is the Lindley--Bartlett effect
\cite{lindley1957,jeffreys1939}: a sharp null beats a diffuse alternative not
on the fit but on the width.
\end{example}

Push $W \to \infty$, the improper limit, and $\Omega \to 0$ at a rate no
finite reading fixes; $B_{10} \to 0$ regardless of $x$. The comparison is no
longer underdetermined, it is ill-posed: the answer depends on a width the
improper prior declines to name \cite{bartlett1957}. Jeffreys' fix is to keep
the prior proper, of fixed scale, so $\Omega$ stays defined \cite{jeffreys1939}.

\paragraph{Coherence is not convergence.}
Now take $n$ readings $x_1, \dots, x_n$ instead of one. Three claims must be kept
apart. The first is \emph{fixed-seed coherence}: with the prior held fixed,
updating one datum at a time equals updating the whole batch, because the
intermediate posterior cancels. This is the machine-checked theorem
\texttt{sequential\_eq\_batch} (Section~\ref{sec:iteration}): with the joint
model fixed, $\cred(\theta \mid x_{1:n})$ is one object whether assembled in $n$
steps or one. That is all it says. It moves no prior and forgets no width; it is
path independence in the data grouping, not a statement about which prior wins.
What it licenses is reading $n$ readings as a single effective measurement of
mean $\bar x$ with noise $1/\sqrt{n}$.

\begin{example}[Accumulating data washes out a broad proper prior]
\label{ex:gravity-washout}
Keep the broad prior $W = 100$ that lost in
Example~\ref{ex:gravity-lindley}, and suppose the true acceleration is
$g = 2$, so $\bar x \to 2$. The effective likelihood ratio for $n$ readings is
$e^{n\bar x^2/2} = e^{2n}$, while the Occam factor on the concentrated
posterior is $\Omega_n \approx \sqrt{2\pi/n}\,/\,W$. So
\[
  B_{10}^{(n)} \;\approx\; e^{2n}\cdot \frac{\sqrt{2\pi/n}}{100}.
\]
At $n = 1$ this is $\approx 0.18$ (gravity loses, as before). At $n = 4$ it is
$e^{8}\cdot\sqrt{\pi/2}/100 \approx 2981 \cdot 0.0125 \approx 37$. At $n = 10$
it is astronomically large. The exponential growth of the likelihood ratio
overruns the fixed width penalty in the denominator. The broad-but-proper prior
that lost on one datum is forgotten once the data accumulate. This is the
\emph{second}, separate claim, convergence under assumptions: the
Bernstein--von Mises behavior of Section~\ref{sec:iteration}, where the posterior
concentrates at rate $1/\sqrt{n}$ and the prior washes out
\cite{vandervaart1998}. The forgetting is the work of the growing likelihood
ratio $e^{2n}$, not of the coherence theorem; the \emph{third} claim, the Occam
factor $\Omega_n$, still divides every term, and prior-volume sensitivity is
removed by no theorem here. Coherence guarantees only that the sequential and
batch routes to this verdict agree.
\end{example}

The contrast with the improper case is the whole point. A proper prior of any
finite width, however badly chosen, is recoverable: enough readings drive
$B_{10}^{(n)}$ to its data-driven verdict, and the coherence theorem guarantees
that the sequential accumulation reaching that verdict is the same as the batch
one. An improper prior, $W = \infty$, has no such limit to recover toward. The
$\Omega = \infty^{-1}$ that made one reading ill-posed makes every reading
ill-posed; no amount of accumulated data names the width, because there was
never a width to name. The seed must be proper for the data to forget it.

This is the same boundary, read on one experiment, that the dictionary of
Section~\ref{sec:correspondence} reads three ways. A finite width is a prior
the evidence can revise (singleton conditioning, $e > 0$). An infinite width is
the unrecoverable case: not the zero of Section~\ref{sec:zero}, but its dual at
the other end, a prior so spread that no finite evidence concentrates it. In
both, the honest report under an unfixed seed is not a number but the set of
verdicts the admissible seeds allow.

\bibliographystyle{plain}
\bibliography{../cred}

\end{document}
